% Inactive reconstruction: not an author input.
% Classical curve factorization, proved through component coverings.
\subsection{The curve of fiber components}

\begin{proposition}\label{draft649:curve-component-factor}
Let $p:X\to C$ be a surjective holomorphic map, where $X$ is a
connected compact complex manifold of dimension $n\geq1$ and
$C$ is a connected compact Riemann surface. There is a
factorization
\[
 X\xrightarrow{q}C'\xrightarrow{r}C,\qquad p=r\circ q,
\]
where $C'$ is a connected compact Riemann surface, $q$ is
surjective and holomorphic with connected fibers, and $r$ is a
finite holomorphic map. The points of $r^{-1}(s)$ correspond to
the connected components of $p^{-1}(s)$, which are the fibers
of $q$ over those points. The factorization is unique up to a
unique isomorphism commuting with $q$ and $r$.

If $t\in r^{-1}(s)$ has ramification index $e_t$, then, as
effective divisors with multiplicities,
\[
 [p^{-1}(s)]=\sum_{t\in r^{-1}(s)}e_t[q^{-1}(t)].
\]
The degree of $r$ is the number of components of a regular
fiber of $p$.
\end{proposition}
\begin{proof}
We first construct the covering over the regular values,
complete it across the critical values, and then extend its
component quotient to $X$.

The critical values form a finite set $D\subset C$, by the
critical-set argument in the proof of
\cref{curve-map:potential-descent}; that argument does not use
connectedness of the fibers. Over $C^\circ=C\setminus D$, the
map $p$ is a proper submersion. The local construction in
\cref{lem:submersion-component-cover} applies there: horizontal
transport over a disk has compact inverse image over its
closure. Compact fibers have finitely many components, and
the resulting local labels are glued by locally constant
permutations. We obtain
\[
 X^\circ\xrightarrow{q^\circ}C'^\circ
             \xrightarrow{r^\circ}C^\circ,
 \qquad X^\circ=p^{-1}(C^\circ),
\]
with $q^\circ$ a holomorphic submersion with connected fibers
and $r^\circ$ a finite unramified covering of degree $d$.
The complement of a proper analytic subset in a connected
complex manifold is connected, so $X^\circ$ and its image
$C'^\circ$ are connected.
Here only the fiber divisor is removed. Locally, the complement
of the zero set of a nonzero holomorphic function in a ball is
path connected: join two outside points through an outside
intermediate point by two complex-line disks. Each restriction
is nonzero because it contains an outside endpoint, and its
isolated zeros can be avoided. Applying this local argument
along a chain of balls proves the required global connectedness.

For each $s\in D$, take a small coordinate disk centered at
$s$ and disjoint from the other critical values. A connected
component of the covering over its puncture is a finite
connected covering of a punctured disk, of some degree $e$.
It has a coordinate $w$ in which $r^\circ(w)=w^e$: the
universal cover of the punctured disk is a half-plane with
projection $z=\exp\zeta$, and the index-$e$ subgroup of its
deck group is generated by $\zeta\mapsto\zeta+2\pi\mathrm i e$.
The coordinate $w=\exp(\zeta/e)$ gives the asserted form.
Attach one point $w=0$ to each such component. There are
finitely many, with their degrees summing to $d$. Gluing these
disks to $C'^\circ$ gives a connected Riemann surface $C'$ and
a finite holomorphic map $r:C'\to C$. The new points have
disjoint disk neighborhoods, so this gluing is Hausdorff.
The inverse image of a small closed base disk is a finite
union of closed disks; over the complement of their interiors,
the original finite covering is proper. Thus $C'$ is compact.

We extend $q^\circ$. Fix $x\in p^{-1}(s)$, $s\in D$. On a
small coordinate ball $B$ about $x$, the image of $p$ lies in
the chosen disk. The complement $B\setminus p^{-1}(s)$ is
connected. Its image under $q^\circ$ therefore lies in a single
punctured disk just attached to $C'$, since those disks are
the connected components of $r^{-1}$ of the punctured base
disk. In its coordinate $w$, the function $w\circ q^\circ$ is
holomorphic and bounded, and
\[
 (w\circ q^\circ)^e=p-s.
\]
The bounded holomorphic removable-singularity theorem extends
it across $p^{-1}(s)\cap B$. To see this without assuming that
the divisor is smooth, first extend at its smooth points by
the one-variable theorem in transverse coordinates. The
singular locus of the reduced divisor has complex codimension
at least two in the ambient manifold, so Hartogs extension
across that locus completes the extension. In dimension one
there are only isolated points. Its $e$th power remains $p-s$ by
the identity theorem, and hence its value is zero on that
fiber. This defines a holomorphic map from $B$ to the attached
disk. The local extensions agree on their dense common
regular locus and therefore glue to $q:X\to C'$ with
$p=r\circ q$. Its image is closed, by compactness, and contains
the dense set $C'^\circ$, so $q$ is surjective.

Every fiber of $q$ is connected. Otherwise write one fiber
as a union of two disjoint nonempty compact subsets and put
them in disjoint open neighborhoods $U_1,U_2$. Properness
gives a disk $V$ about their base point $t$ such that
\[
 q^{-1}(V)\subset U_1\cup U_2.
\]
The map $q$ is open: at any point, a complex line on which
the nonconstant germ $q-t$ has its first nonzero term gives
an open one-variable restriction. Consequently both
$q(U_i)$ contain a neighborhood of $t$. A nearby point of
$C'^\circ$ then has a fiber meeting both neighborhoods,
contradicting connectedness of the regular fibers.

The fibers of $p$ are the disjoint union of these connected
$q$-fibers over the finitely many points of $r^{-1}(s)$.
They are its components. Locally at $t$, the formula for a
finite map is $r(w)-s=w^{e_t}$ after changing the coordinate
by a holomorphic root of a nonvanishing factor. Pulling back
this equality by $q$ gives the divisor multiplicities stated
above. Over $C^\circ$, $r$ is unramified and each sheet is
one component, proving the degree assertion.

Finally, any factorization with the stated properties has
over $C^\circ$ precisely this covering of fiber components.
It is unramified there: a critical point of its finite map,
together with surjectivity of its first factor, would make
$p$ critical over that base point. Thus the component
identification is holomorphic on $C'^\circ$. On every
punctured disk it identifies two completions of the same
finite covering, and the coordinates $w^e$ above extend it
biholomorphically across the added point. It gives a global
isomorphism; density and surjectivity make it unique and
give commutation with both maps.
\end{proof}

\begin{corollary}\label{draft649:full-curve-model-factor}
Let $p:X\to C$ be any surjective holomorphic map from a
connected compact complex manifold of dimension $n\geq2$
onto a connected compact Riemann surface. Let $\beta$ be a
smooth real form on $C$ and $A=\int_C\beta$. Then
$\PSH(X,p^*\beta)$ is nonempty if and only if $A\geq0$.
For $A\geq0$, use the factorization above and put
$\sigma=r^*\beta$, whose area is $dA$. The full ordinary
model space is the pullback under $q$ of the curve model
space in background $\sigma$. Every potential has zero
extended top-degree measure and every ordinary type has
zero upper and lower singularity volumes.

The upstairs polar-trace metric normalized by the original
$p$-fiber volume is complete, geodesic, and independent of
the Gauduchon form, with diameter $2A$. Its analytic closure
is the ideal subspace, and ideal projection is its unique
nearest-point, $1$-Lipschitz retraction. The curve-factor
normalization gives $d$ times this metric. At $A=0$ the
model space is a single smooth model.
\end{corollary}
\begin{proof}
The proposition supplies the factorization assumed in
\cref{fibration:finite-curve-factor}. Since
$\int_{C'}r^*\beta=dA$, the nonemptiness and zero-mass
assertions follow from
\cref{fibration:curve-envelope-pullback} applied to $q$.
The finite-factor metric theorem gives the remaining
conclusions; multiplying a metric by the positive constant
$d^{-1}$ preserves completeness, geodesics, analytic closure,
nearest-point projection and its Lipschitz constant.
\end{proof}
